% =============================================================================
% APPENDIX VB: METHOD-PAPERQA: Paper Integration Quality Assurance
% =============================================================================
% Module: BCM2_METHOD_QA
% Version: 1.0 (February 2026)
% Status: Final
%
% Defines the quality assurance workflow for paper integration, including
% the 13-component validation checklist, compliance thresholds, and the
% 6-Factor Framework for CORE integration decisions.
%
% RELATIONSHIP TO BM: Appendix BM defines the 2D Classification System
% (Content Level x Integration Level). This appendix (VB) provides the
% operational QA procedures for ensuring integration quality.
%
% =============================================================================

% =============================================================================
% CROSS-REFERENCE STRUCTURE
% =============================================================================
%
% CHAPTER LINKAGE (Required):
% - Primary Chapter: CLAUDE.md (Framework Documentation)
% - Secondary Chapters: Quality Assurance Workflows
%
% DEPENDENCIES (what this appendix REQUIRES):
% - Appendix UNMAPPED_BBB (CORE-WHERE): Parameter values
% - Appendix UNMAPPED_MS (REF-MODELSPACE): Theory catalog
% - Data: paper-references/*.yaml, case-registry.yaml
%
% DEPENDENTS (what REQUIRES this appendix):
% - All paper integration workflows
% - Quality assurance processes
% - LIT-* appendices
%
% =============================================================================

\begin{tcolorbox}[colback=purple!5!white, colframe=purple!75!black,
    title=Chapter Linkage]

\textbf{Primary Reference:} CLAUDE.md Section ``/integrate-paper''

\textbf{Related Workflows:}
\begin{itemize}[nosep]
\item Evidence Integration Pipeline (EIP)
\item Level 5 Paper Integration Workflow
\item 6-Factor Framework for CORE Integration
\end{itemize}

\textbf{Validation Script:} \texttt{scripts/validate\_level5\_integration.py}

\end{tcolorbox}

\chapter{METHOD-PAPERQA: Paper Integration Quality Assurance}
\label{app:paperqa}


\begin{tcolorbox}[colback=blue!5!white, colframe=blue!75!black,
    title=Quick Reference: Appendix UNMAPPED_VB]

\textbf{Appendix:} VB (METHOD-PAPERQA)

\textbf{Core Question:} How do we ensure quality in paper integration?

\textbf{Companion Appendix:} BM (METHOD-PAPERINT) defines 2D Classification; this appendix provides QA procedures.

\textbf{Key Concepts:}
\begin{itemize}[nosep]
\item 13-Component Validation System
\item 5 Integration Levels (MINIMAL $\rightarrow$ FULL)
\item ATOMIC ID RULE (Paper-YAML as SSOT)
\item 6-Factor Framework for CORE Integration
\end{itemize}

\textbf{Cross-References:} BBB (WHERE), MS (Theory Catalog), LIT-* Appendices
\end{tcolorbox}

\begin{abstract}
This appendix defines the standardized Paper Integration Workflow for the
Evidence-Based Framework (EBF). It specifies five integration levels
(MINIMAL to FULL), the 13-component validation system, compliance thresholds,
and the 6-Factor Framework for deciding CORE appendix integration. The workflow
ensures consistent, high-quality paper integration with full traceability.
\end{abstract}


\begin{tcolorbox}[
    colback=orange!5!white,
    colframe=orange!75!black,
    title={\textbf{Appendix Scope: Ziel / In-Scope / Out-of-Scope / Constraints}},
    fonttitle=\bfseries
]

\textbf{Ziel (Objective):}
\begin{itemize}[nosep]
    \item Define standardized workflow for paper integration at all quality levels
\end{itemize}

\textbf{In-Scope:}
\begin{itemize}[nosep]
    \item 5 Integration Levels (MINIMAL, STANDARD, CASE, THEORY, FULL)
    \item 13-Component Validation Checklist
    \item Compliance Thresholds per Level
    \item 6-Factor Framework for CORE Integration Decisions
    \item ATOMIC ID RULE (Paper-YAML as Single Source of Truth)
\end{itemize}

\textbf{Out-of-Scope (delegiert an andere Appendices):}
\begin{itemize}[nosep]
    \item Parameter values $\rightarrow$ Appendix UNMAPPED_BBB (CORE-WHERE)
    \item Theory definitions $\rightarrow$ Appendix UNMAPPED_MS (REF-MODELSPACE)
    \item Case structures $\rightarrow$ Case Registry (data/case-registry.yaml)
\end{itemize}

\textbf{Constraints (Anwendungsgrenzen):}
\begin{itemize}[nosep]
    \item Requires existing Paper-YAML infrastructure
    \item Compliance thresholds are minimum requirements, not targets
    \item 6-Factor Framework applies only to Level 4+ integrations
\end{itemize}

\textbf{Lieferobjekte (Deliverables):}
\begin{itemize}[nosep]
    \item 13-Component Checklist (Table~\ref{tab:13-components})
    \item Compliance Threshold Matrix (Table~\ref{tab:compliance-thresholds})
    \item 6-Factor Decision Framework (Table~\ref{tab:6-factor})
    \item Worked Example: Bénabou et al. (2024)
\end{itemize}

\end{tcolorbox}


%%%%%%%%%%%%%%%%%%%%%%%%%%%%%%%%%%%%%%%%%%%%%%%%%%%%%%%%%%%%%%%%%%%%%%%%%%%%%%%
\section{The Fundamental Question}
\label{sec:paperint-fundamental}
%%%%%%%%%%%%%%%%%%%%%%%%%%%%%%%%%%%%%%%%%%%%%%%%%%%%%%%%%%%%%%%%%%%%%%%%%%%%%%%

\subsection{Why This Matters}

The EBF framework integrates 2,500+ scientific papers. Without a standardized
integration workflow, paper quality varies, cross-references become inconsistent,
and the framework loses its scientific rigor.

\subsection{The Core Problem}

Different papers require different integration depths. A working paper with
a single parameter estimate requires minimal integration, while a foundational
paper introducing a new theory category requires comprehensive integration
across all EBF components.

\begin{tcolorbox}[colback=red!5!white, colframe=red!75!black,
    title=The Fundamental Question]
\textbf{How do we ensure consistent, traceable, and appropriate-depth integration
of scientific papers into the EBF framework?}
\end{tcolorbox}


%%%%%%%%%%%%%%%%%%%%%%%%%%%%%%%%%%%%%%%%%%%%%%%%%%%%%%%%%%%%%%%%%%%%%%%%%%%%%%%
\section{The Five Integration Levels}
\label{sec:paperint-levels}
%%%%%%%%%%%%%%%%%%%%%%%%%%%%%%%%%%%%%%%%%%%%%%%%%%%%%%%%%%%%%%%%%%%%%%%%%%%%%%%

\begin{table}[htbp]
\centering
\caption{Paper Integration Levels}
\label{tab:integration-levels}
\renewcommand{\arraystretch}{1.3}
\begin{tabular}{@{}clcl@{}}
\toprule
\textbf{Level} & \textbf{Name} & \textbf{Components} & \textbf{Typical Use Case} \\
\midrule
1 & MINIMAL & 1-2 & Working paper, single citation \\
2 & STANDARD & 3-5 & Peer-reviewed with theory support \\
3 & CASE & 6-8 & Field study with case registry entry \\
4 & THEORY & 9-11 & Paper extending existing theory \\
5 & FULL & 12-13 & Foundational paper, new domain \\
\bottomrule
\end{tabular}
\end{table}

\subsection{Level 1: MINIMAL}

\textbf{Components:} BibTeX entry only.

\textbf{Time:} $\sim$5 minutes.

\textbf{Use when:} Paper provides single citation, no parameter extraction needed.

\subsection{Level 2: STANDARD}

\textbf{Components:} BibTeX + theory\_support + use\_for fields.

\textbf{Time:} $\sim$15 minutes.

\textbf{Use when:} Peer-reviewed paper supports existing EBF theories.

\subsection{Level 3: CASE}

\textbf{Components:} Level 2 + Case Registry entry.

\textbf{Time:} $\sim$25 minutes.

\textbf{Use when:} Paper contains worked example suitable for case registry.

\subsection{Level 4: THEORY}

\textbf{Components:} Level 3 + Paper-YAML + Theory Catalog entry + LIT-Appendix.

\textbf{Time:} $\sim$45 minutes.

\textbf{Use when:} Paper extends or refines existing EBF theory.

\subsection{Level 5: FULL}

\textbf{Components:} All 13 components including CORE appendix integration.

\textbf{Time:} $\sim$90 minutes.

\textbf{Use when:} Foundational paper introducing new concepts, domains, or paradigms.


%%%%%%%%%%%%%%%%%%%%%%%%%%%%%%%%%%%%%%%%%%%%%%%%%%%%%%%%%%%%%%%%%%%%%%%%%%%%%%%
\section{The 13-Component Validation System}
\label{sec:paperint-components}
%%%%%%%%%%%%%%%%%%%%%%%%%%%%%%%%%%%%%%%%%%%%%%%%%%%%%%%%%%%%%%%%%%%%%%%%%%%%%%%

\begin{tcolorbox}[colback=gray!5!white, colframe=gray!75!black,
    title=Axiom PINT-1: Component Completeness]
\textbf{Statement:} A paper at Level $L$ is complete if and only if all
components required for that level are present and validated.

\textbf{Interpretation:} Partial integration creates inconsistencies;
completeness at each level ensures traceability.

\textbf{Implication:} Papers should not be marked as integrated until
all level-appropriate components are verified.
\end{tcolorbox}

\begin{table}[htbp]
\centering
\caption{13-Component Validation Checklist}
\label{tab:13-components}
\renewcommand{\arraystretch}{1.3}
\begin{tabular}{@{}clccccc@{}}
\toprule
\textbf{\#} & \textbf{Component} & \textbf{L1} & \textbf{L2} & \textbf{L3} & \textbf{L4} & \textbf{L5} \\
\midrule
1 & Paper-YAML exists & -- & -- & -- & $\star$ & $\star$ \\
2 & Full text archived & -- & -- & -- & $\star$ & $\star$ \\
3 & Content level $\geq$ L1 & -- & $\circ$ & $\circ$ & $\star$ & $\star$ \\
4 & BibTeX entry & $\star$ & $\star$ & $\star$ & $\star$ & $\star$ \\
5 & theory\_support field & -- & $\star$ & $\star$ & $\star$ & $\star$ \\
6 & use\_for field & -- & $\star$ & $\star$ & $\star$ & $\star$ \\
7 & Case Registry entry & -- & -- & $\star$ & $\star$ & $\star$ \\
8 & Chapter-Appendix mapping & -- & -- & -- & $\star$ & $\star$ \\
9 & Theory Catalog entry & -- & -- & -- & $\circ$ & $\star$ \\
10 & Parameter Registry & -- & -- & -- & $\circ$ & $\star$ \\
11 & LIT-Appendix integration & -- & -- & -- & $\star$ & $\star$ \\
12 & CORE-Appendix cross-ref & -- & -- & -- & -- & $\star$ \\
13 & 10C dimension mapping & -- & -- & -- & -- & $\star$ \\
\bottomrule
\end{tabular}
\begin{tablenotes}
\small
\item $\star$ = Required, $\circ$ = Optional, -- = Not applicable
\end{tablenotes}
\end{table}


%%%%%%%%%%%%%%%%%%%%%%%%%%%%%%%%%%%%%%%%%%%%%%%%%%%%%%%%%%%%%%%%%%%%%%%%%%%%%%%
\section{Compliance Thresholds}
\label{sec:paperint-compliance}
%%%%%%%%%%%%%%%%%%%%%%%%%%%%%%%%%%%%%%%%%%%%%%%%%%%%%%%%%%%%%%%%%%%%%%%%%%%%%%%

\begin{tcolorbox}[colback=gray!5!white, colframe=gray!75!black,
    title=Axiom PINT-2: Threshold Enforcement]
\textbf{Statement:} Paper integration at Level $L$ requires compliance score
$C(L) \geq T(L)$ where $T(L)$ is the threshold for level $L$.

\textbf{Interpretation:} Higher integration levels require stricter compliance.

\textbf{Implication:} Papers below threshold must complete missing components
before being marked as integrated at that level.
\end{tcolorbox}

\begin{table}[htbp]
\centering
\caption{Compliance Thresholds by Integration Level}
\label{tab:compliance-thresholds}
\renewcommand{\arraystretch}{1.3}
\begin{tabular}{@{}clcc@{}}
\toprule
\textbf{Level} & \textbf{Name} & \textbf{Threshold} & \textbf{Interpretation} \\
\midrule
3 & CASE & $\geq 60\%$ & Minimum for case-level integration \\
4 & THEORY & $\geq 75\%$ & Required for theory extensions \\
5 & FULL & $\geq 85\%$ & Required for foundational papers \\
\bottomrule
\end{tabular}
\end{table}

\textbf{Validation Command:}
\begin{verbatim}
python scripts/validate_level5_integration.py PAP-{bibtex_key}
\end{verbatim}


%%%%%%%%%%%%%%%%%%%%%%%%%%%%%%%%%%%%%%%%%%%%%%%%%%%%%%%%%%%%%%%%%%%%%%%%%%%%%%%
\section{The ATOMIC ID RULE}
\label{sec:paperint-atomic}
%%%%%%%%%%%%%%%%%%%%%%%%%%%%%%%%%%%%%%%%%%%%%%%%%%%%%%%%%%%%%%%%%%%%%%%%%%%%%%%

\begin{tcolorbox}[colback=gray!5!white, colframe=gray!75!black,
    title=Axiom PINT-3: ATOMIC ID RULE]
\textbf{Statement:} The Paper-YAML file (\texttt{data/paper-references/PAP-\{key\}.yaml})
is the Single Source of Truth (SSOT) for all paper metadata.

\textbf{Interpretation:} All other references (BibTeX, LIT-Appendix, Case Registry)
must derive from and be consistent with the Paper-YAML.

\textbf{Implication:} When updating paper information, always update the
Paper-YAML first, then propagate changes to dependent files.
\end{tcolorbox}

\subsection{ID Format}

\begin{equation}
\text{Paper-ID} = \texttt{PAP-}\{\text{firstauthor}\}\_\{\text{year}\}\_\{\text{shortname}\}
\end{equation}

\textbf{Examples:}
\begin{itemize}[nosep]
\item \texttt{PAP-benabou\_2024\_ends\_means}
\item \texttt{PAP-fehr\_1999\_inequity}
\item \texttt{PAP-kahneman\_1979\_prospect}
\end{itemize}

\subsection{SSOT Hierarchy}

\begin{enumerate}
\item \textbf{Paper-YAML} (authoritative) $\rightarrow$ All metadata
\item \textbf{Paper-Text} (archival) $\rightarrow$ Full text in Markdown
\item \textbf{BibTeX} (derived) $\rightarrow$ Citation data only
\item \textbf{Case Registry} (derived) $\rightarrow$ Links via paper\_id
\item \textbf{LIT-Appendix} (derived) $\rightarrow$ References Paper-YAML
\end{enumerate}


%%%%%%%%%%%%%%%%%%%%%%%%%%%%%%%%%%%%%%%%%%%%%%%%%%%%%%%%%%%%%%%%%%%%%%%%%%%%%%%
\section{The 6-Factor Framework for CORE Integration}
\label{sec:paperint-6factor}
%%%%%%%%%%%%%%%%%%%%%%%%%%%%%%%%%%%%%%%%%%%%%%%%%%%%%%%%%%%%%%%%%%%%%%%%%%%%%%%

\begin{tcolorbox}[colback=gray!5!white, colframe=gray!75!black,
    title=Axiom PINT-4: CORE Integration Decision]
\textbf{Statement:} A paper should extend a CORE appendix if and only if
at least 3 of the 6 factors evaluate to ``Yes'' AND factor F1 (Structural Novelty)
evaluates to ``Yes''.

\textbf{Interpretation:} CORE appendices should only be modified when
a paper introduces genuinely new structural components, not mere applications.

\textbf{Implication:} Most papers (even Level 5) will have cross-references
to CORE appendices but will not modify them.
\end{tcolorbox}

\begin{table}[htbp]
\centering
\caption{6-Factor Framework for CORE Integration}
\label{tab:6-factor}
\renewcommand{\arraystretch}{1.3}
\begin{tabular}{@{}clp{5cm}cc@{}}
\toprule
\textbf{\#} & \textbf{Factor} & \textbf{Question} & \textbf{Yes} & \textbf{No} \\
\midrule
F1 & Structural Novelty & Does paper introduce NEW component/dimension? & Extend & Reference \\
F2 & Axiom Contribution & Does paper enable NEW axiom? & Formulate & Support \\
F3 & Dimensionality & Does CORE-space dimension change? & Extend & Same \\
F4 & Parameter vs Structure & Is contribution structural or parametric? & CORE & WHERE \\
F5 & Universality & Does contribution apply GENERALLY? & CORE & LIT/DOMAIN \\
F6 & Mechanism & Does paper explain NEW mechanism? & CORE & Application \\
\bottomrule
\end{tabular}
\end{table}

\subsection{Decision Tree}

\begin{verbatim}
Paper P to integrate into CORE-X?
│
├── F1: New structural component for X?
│   ├── NO → Cross-reference only
│   └── YES ──┐
│             │
│             ├── F5: Applies universally (not domain-specific)?
│             │   ├── NO → LIT/DOMAIN appendix, cross-ref to CORE
│             │   └── YES ──┐
│             │             │
│             │             └── Extend CORE-X with:
│             │                 • New component/dimension
│             │                 • New axiom (if F2=YES)
│             │                 • Cross-reference to LIT for details
│             │
│             └── Parameter values → WHERE (BBB), not CORE-X
\end{verbatim}


%%%%%%%%%%%%%%%%%%%%%%%%%%%%%%%%%%%%%%%%%%%%%%%%%%%%%%%%%%%%%%%%%%%%%%%%%%%%%%%
\section{Worked Example: Bénabou, Falk \& Henkel (2024)}
\label{sec:paperint-example}
%%%%%%%%%%%%%%%%%%%%%%%%%%%%%%%%%%%%%%%%%%%%%%%%%%%%%%%%%%%%%%%%%%%%%%%%%%%%%%%

\subsection{Paper Overview}

\textbf{Title:} ``Ends versus Means: Kantians, Utilitarians, and Moral Decisions''

\textbf{Key Finding:} No stable moral types---cross-context correlation $r_{\text{EVM}} \approx 0$.

\textbf{Target Level:} Level 4 (THEORY)

\subsection{Component Completion Progress}

\begin{table}[htbp]
\centering
\caption{Integration Progress: PAP-benabou\_2024\_ends\_means}
\renewcommand{\arraystretch}{1.2}
\begin{tabular}{@{}clcc@{}}
\toprule
\textbf{\#} & \textbf{Component} & \textbf{Initial} & \textbf{Final} \\
\midrule
1 & Paper-YAML exists & \checkmark & \checkmark \\
2 & Full text archived & \checkmark & \checkmark \\
3 & Content level $\geq$ L1 & \checkmark & \checkmark \\
4 & BibTeX entry & \checkmark & \checkmark \\
5 & theory\_support field & \checkmark & \checkmark \\
6 & use\_for field & \checkmark & \checkmark \\
7 & Case Registry entry & $\times$ & \checkmark \\
8 & Chapter-Appendix mapping & $\times$ & \checkmark \\
9 & Theory Catalog entry & \checkmark & \checkmark \\
10 & Parameter Registry & \checkmark & \checkmark \\
11 & LIT-Appendix integration & $\times$ & \checkmark \\
\midrule
& \textbf{Compliance Score} & \textbf{15.4\%} & \textbf{84.6\%} \\
\bottomrule
\end{tabular}
\end{table}

\subsection{6-Factor Analysis}

\begin{table}[htbp]
\centering
\caption{6-Factor Analysis: Bénabou et al. (2024)}
\renewcommand{\arraystretch}{1.2}
\begin{tabular}{@{}clcl@{}}
\toprule
\textbf{\#} & \textbf{Factor} & \textbf{Value} & \textbf{Reasoning} \\
\midrule
F1 & Structural Novelty & NO & Uses existing $\theta = f(\Psi)$ framework \\
F2 & Axiom Contribution & NO & Supports existing PCT axioms \\
F3 & Dimensionality & NO & No new dimensions \\
F4 & Parameter vs Structure & Parameter & $r_{\text{EVM}} \approx 0$ is a parameter \\
F5 & Universality & YES & General moral reasoning finding \\
F6 & Mechanism & NO & Confirms context-dependence mechanism \\
\midrule
& \textbf{Decision} & \multicolumn{2}{l}{Cross-reference to CORE, no extension} \\
\bottomrule
\end{tabular}
\end{table}

\textbf{Result:} Paper integrated at Level 4 with:
\begin{itemize}[nosep]
\item Case UNMAPPED_CAS-903 in Case Registry
\item Axiom UNMAPPED_RB-11 in LIT-BENABOU appendix
\item Cross-references to CORE-WHEN (V), CORE-HOW (B), CORE-AWARE (AU)
\item Parameters PAR-MR-001 to PAR-MR-004 in Parameter Registry
\end{itemize}


%%%%%%%%%%%%%%%%%%%%%%%%%%%%%%%%%%%%%%%%%%%%%%%%%%%%%%%%%%%%%%%%%%%%%%%%%%%%%%%
\section{Quality Assurance Integration}
\label{sec:paperint-qa}
%%%%%%%%%%%%%%%%%%%%%%%%%%%%%%%%%%%%%%%%%%%%%%%%%%%%%%%%%%%%%%%%%%%%%%%%%%%%%%%

\begin{tcolorbox}[colback=green!5!white, colframe=green!75!black,
    title=Practical Implications]
\begin{enumerate}
\item \textbf{Always validate before marking complete:} Run
      \texttt{validate\_level5\_integration.py} for Level 4+ papers.
\item \textbf{Follow PDCA cycle:} Plan components, Do integration,
      Check compliance, Act on gaps.
\item \textbf{Document in Lessons Learned:} Record workflow insights
      in \texttt{quality/lessons\_learned.md}.
\item \textbf{Update checklist:} Track Level 4+ integrations in
      \texttt{quality/checklist.md}.
\end{enumerate}
\end{tcolorbox}

\subsection{PDCA Cycle for Paper Integration}

\begin{enumerate}
\item \textbf{Plan:} Determine target level, identify required components
\item \textbf{Do:} Complete each component in order
\item \textbf{Check:} Run validation script, verify compliance $\geq$ threshold
\item \textbf{Act:} Fix gaps, document lessons learned
\end{enumerate}


%%%%%%%%%%%%%%%%%%%%%%%%%%%%%%%%%%%%%%%%%%%%%%%%%%%%%%%%%%%%%%%%%%%%%%%%%%%%%%%
\section{Summary}
\label{sec:paperint-summary}
%%%%%%%%%%%%%%%%%%%%%%%%%%%%%%%%%%%%%%%%%%%%%%%%%%%%%%%%%%%%%%%%%%%%%%%%%%%%%%%

\begin{tcolorbox}[colback=green!10!white, colframe=green!75!black,
    title=\textbf{Appendix UNMAPPED_VB: Summary}]

\textbf{Core Question:} How do we systematically integrate papers into EBF?

\textbf{Main Contribution:}
\begin{itemize}[nosep]
    \item 5-level integration system (MINIMAL $\rightarrow$ FULL)
    \item 13-component validation checklist
    \item 6-Factor Framework for CORE integration decisions
    \item ATOMIC ID RULE for traceability
\end{itemize}

\textbf{Key Outputs:}
\begin{itemize}[nosep]
    \item Compliance thresholds: L3 $\geq 60\%$, L4 $\geq 75\%$, L5 $\geq 85\%$
    \item Validation script: \texttt{validate\_level5\_integration.py}
\end{itemize}

\textbf{Integration:} This appendix connects to BBB (WHERE) for parameters,
MS (Theory Catalog) for theories, and all LIT-* appendices for literature.

\end{tcolorbox}


%%%%%%%%%%%%%%%%%%%%%%%%%%%%%%%%%%%%%%%%%%%%%%%%%%%%%%%%%%%%%%%%%%%%%%%%%%%%%%%
\section{Glossary of Symbols}
\label{sec:paperint-glossary}
%%%%%%%%%%%%%%%%%%%%%%%%%%%%%%%%%%%%%%%%%%%%%%%%%%%%%%%%%%%%%%%%%%%%%%%%%%%%%%%

\begin{tcolorbox}[colback=blue!5!white, colframe=blue!75!black]
\textbf{Central Reference:} For complete symbol definitions, see
\textbf{Appendix G} (Glossary).
\end{tcolorbox}

\begin{table}[htbp]
\centering
\caption{Symbols Introduced in Appendix UNMAPPED_VB}
\label{tab:paperint-symbols}
\renewcommand{\arraystretch}{1.3}
\begin{tabular}{@{}clp{6cm}c@{}}
\toprule
\textbf{Symbol} & \textbf{Name} & \textbf{Definition} & \textbf{Status} \\
\midrule
$L$ & Integration Level & Paper integration depth (1--5) & NEW \\
$C(L)$ & Compliance Score & Percentage of components completed & NEW \\
$T(L)$ & Threshold & Minimum compliance for level $L$ & NEW \\
$r_{\text{EVM}}$ & Cross-EVM Correlation & Avg. correlation across moral contexts & $\checkmark$ \\
$\theta$ & Deontological Weight & Weight on rule-following vs. outcomes & $\checkmark$ \\
$\Psi$ & Context Vector & 8-dimensional context specification & $\checkmark$ \\
\bottomrule
\end{tabular}
\end{table}


%%%%%%%%%%%%%%%%%%%%%%%%%%%%%%%%%%%%%%%%%%%%%%%%%%%%%%%%%%%%%%%%%%%%%%%%%%%%%%%
\section{Open Issues and Research Agenda}
\label{sec:paperint-open}
%%%%%%%%%%%%%%%%%%%%%%%%%%%%%%%%%%%%%%%%%%%%%%%%%%%%%%%%%%%%%%%%%%%%%%%%%%%%%%%

\begin{table}[htbp]
\centering
\caption{Research Roadmap for Paper Integration}
\label{tab:paperint-roadmap}
\renewcommand{\arraystretch}{1.3}
\begin{tabular}{@{}clcl@{}}
\toprule
\textbf{\#} & \textbf{Issue} & \textbf{Priority} & \textbf{Status} \\
\midrule
1 & Automated compliance scoring & HIGH & Implemented \\
2 & Pre-commit hook integration & HIGH & Implemented \\
3 & Paper queue management & MEDIUM & Implemented \\
4 & Cross-database consistency checks & MEDIUM & Planned \\
5 & Automated LIT-appendix sync & LOW & Partial \\
\bottomrule
\end{tabular}
\end{table}


%%%%%%%%%%%%%%%%%%%%%%%%%%%%%%%%%%%%%%%%%%%%%%%%%%%%%%%%%%%%%%%%%%%%%%%%%%%%%%%
\section{References}
\label{sec:paperint-references}
%%%%%%%%%%%%%%%%%%%%%%%%%%%%%%%%%%%%%%%%%%%%%%%%%%%%%%%%%%%%%%%%%%%%%%%%%%%%%%%

\begin{tcolorbox}[colback=blue!5!white, colframe=blue!75!black]
\textbf{Master Bibliography:} See \texttt{bibliography/bcm\_master.bib}
\end{tcolorbox}

\subsection{Key References for This Appendix}

\begin{itemize}[nosep]
    \item \textbf{Paper Integration:} CLAUDE.md Section ``/integrate-paper''
    \item \textbf{Quality Assurance:} \texttt{quality/checklist.md}
    \item \textbf{Lessons Learned:} \texttt{quality/lessons\_learned.md}
\end{itemize}

\nocite{bcm_master}

\vspace{1em}
\hrule
\vspace{0.5em}

\noindent\textit{Cross-references:}
Appendix G (Glossary),
Appendix UNMAPPED_BBB (CORE-WHERE),
Appendix UNMAPPED_MS (Theory Catalog),
Appendix UNMAPPED_RB (LIT-BENABOU).

\noindent\textit{Master Bibliography:} \texttt{bibliography/bcm\_master.bib}

% =============================================================================
% END OF APPENDIX VB
% =============================================================================
